% Chapter 2 entries ECHO .. REN (AN-01).

% ECHO -- os/milton/command.c builtin_echo (no arg: "ECHO is on"; else the
% text as typed).  ECHO ON/OFF in batch files: Chapter 4 (batch.c).
% Gates: test-command, test-batch.
\begin{DosEntry}{ECHO}{}{internal}
\Purpose Displays a message, or the state of batch-file echoing.
\Format \SyntaxLine{\kw{ECHO} \opt{\ph{message}}}
\Type{Internal}
\Remarks
\Param{message} is the text to display.  It is displayed exactly as you
type it, in capital and small letters.

Typed with no message, \cmd{ECHO} displays \texttt{ECHO is on}.

In a batch file, \cmd{ECHO ON} and \cmd{ECHO OFF} control whether the
command lines of the batch file are displayed as they are carried out.  See
``ECHO Subcommand'' in Chapter~4.
\Example
\Shot{echo}
\end{DosEntry}

% EXIT -- os/milton/command.c dispatch_line CMD_EXIT (a permanent processor
% ignores EXIT without a message; a secondary returns to its parent).
% Gates: test-shell-permanent, test-exit-handles.
\begin{DosEntry}{EXIT}{}{internal}
\Purpose Ends a secondary command processor and returns to the program that
started it.
\Format \SyntaxLine{\kw{EXIT}}
\Type{Internal}
\Remarks
A program can start a second copy of the command processor so that you can
type commands without leaving the program.  \cmd{EXIT} ends that copy and
returns you to the program.

The command processor that InitechDOS starts when the system is turned on
is permanent.  At its prompt \cmd{EXIT} has no effect and displays nothing.
\Example
At the \texttt{A:\bs>} prompt of a freshly started system, type
\type{EXIT} and press \key{Enter}.  The prompt is displayed again, because
the command processor started with the system is permanent.
\end{DosEntry}

% LABEL -- os/milton/dosutil.c label_main / label_read / label_write
% (current drive only; prompts; the invalid characters; 11 characters,
% capitals; root VOLLABEL entry + extended BPB label field; messages).
% Shipped as APPS\LABEL.COM.  Gates: test-label-vol, test-dosutil-emu.
\begin{DosEntry}{LABEL}{Volume Label}{external}
\Purpose Creates, changes or deletes the volume label of a disk.
\Format \SyntaxLine{\opt{\ph{d}\kw{:}}\opt{\ph{path}}\kw{LABEL}
  \opt{\ph{d}\kw{:}}\opt{\ph{label}}}
\Type{External}
\Remarks
\Param{d:} is the drive of the disk.  Only the current drive can be
labeled; another drive letter displays \texttt{Invalid drive
specification}.

\Param{label} is the new volume label, up to 11 characters.  Letters are
changed to capitals; characters after the eleventh are ignored.  A label
may contain spaces, but not the characters
\texttt{* ? / \textbackslash{} | . , ; : + = < > [ ] ( ) \& \textasciicircum{} "};
\cmd{LABEL} displays \texttt{Invalid characters in volume label} and
leaves the label unchanged.

If you type \cmd{LABEL} with no label, it displays the current label (or
\texttt{has no label}) and asks for a new one:
\begin{TypeIn}
Volume label (11 characters, ENTER for none)?
\end{TypeIn}
Type the new label and press \key{Enter}.  If you press \key{Enter} alone
and the disk has a label, \cmd{LABEL} asks
\texttt{Delete current volume label (Y/N)?}  Type \type{Y} to delete the
label; any other answer keeps it.

The label is shown by \cmd{DIR}, \cmd{VOL}, \cmd{CHKDSK} and \cmd{TREE}.
If the root directory has no room for the label, \cmd{LABEL} displays
\texttt{No room in root directory}.
\Example
To give the practice diskette the label \file{REPORTS}, answering the
question, and display it:
\Shot{label}
To set the label in one command, and an attempt with a period in it:
\Shot{labelset}
\end{DosEntry}

% MD -- os/milton/command.c builtin_md (MSG-DOS-0011, MSG-DOS-0017; silent
% success); int21.c AH=39h.  Directory date stamp: initech-dwbl.
% Gates: test-ut6d, test-fat12-mkdir, test-fat-subdir.
\begin{DosEntry}{MD}{Make Directory}{internal}
\AlsoNamed{MKDIR}
\Purpose Creates a subdirectory.
\Format \SyntaxLine{\kw{MD} \opt{\ph{d}\kw{:}}\ph{path}}
  \SyntaxLine{\kw{MKDIR} \opt{\ph{d}\kw{:}}\ph{path}}
\Type{Internal}
\Remarks
\Param{path} names the new directory.  A path that begins with a backslash
starts at the root directory; any other path starts at the current
directory.  Every directory named before the last one must already exist.
A directory name has the same form as a file name.

\cmd{MD} displays nothing when it succeeds.  A new directory contains two
entries: \file{.}, the directory itself, and \file{..}, its parent.

If the directory cannot be created (for example, because a file or
directory of that name already exists), \cmd{MD} displays
\texttt{Unable to create directory}.  With no path it displays
\texttt{Required parameter missing}.

\draft{initech-dwbl: the shipped build stamps new directories 0-00-80
12:00a (see the example); DOS 3.3 stamps them with the current date and
time}
The new directory is stamped with the current date and time.
\Example
To create the directory \file{REPORTS}, make it the current directory and
list it:
\Shot{md}
\end{DosEntry}

% PATH -- not a command in the shipped build (cmd_classify has no PATH row;
% typing it displays "Bad command or file name").  initech-sy5f.  The PATH
% environment variable itself IS used (command.c run_external,
% cmd_path_candidates; SET PATH= gated by test-command, test-env).
\begin{DosEntry}[initech-sy5f: PATH is not yet an internal command in the
shipped build; the entry describes DOS 3.3.  SET PATH= works today]{PATH}{}{internal}
\Purpose Sets or displays the list of directories searched for commands.
\Format \SyntaxLine{\kw{PATH} \opt{\opt{\ph{d}\kw{:}}\ph{path}\opt{\kw{;}\opt{\ph{d}\kw{:}}\ph{path}}\rep}}
  \SyntaxLine{\kw{PATH} \kw{;}}
\Type{Internal}
\Remarks
\Param{path} names a directory to search.  Separate the directories with
semicolons.  When a command is not an internal command and is not in the
current directory, InitechDOS searches these directories, in order.

Typed with no parameter, \cmd{PATH} displays the current list, as
\texttt{PATH=}\ph{list}.  \type{PATH ;} clears the list, so that only the
current directory is searched.

\cmd{PATH} stores the list in the environment variable \texttt{PATH}; you
can also set it with \type{SET PATH=A:\bs APPS}, as the practice
diskette's \file{AUTOEXEC.BAT} does.
\Example
\type{PATH A:\bs APPS} makes InitechDOS search the directory
\file{\bs APPS} for commands.
\end{DosEntry}

% PROMPT -- os/milton/command.c builtin_prompt (empty -> $P$G; env PROMPT;
% "Out of environment space"), cmd_render_prompt ($P $G $L $B $Q $N $T $D
% $$ $_; any other $x prints x).  Defects found: $D shows 00-00-0000 (the
% REPL passes no date) and $$ cuts the prompt (AH=09h '$' terminator) --
% beads needed, WL-0119.  Gates: test-command, test-shell.
\begin{DosEntry}{PROMPT}{}{internal}
\Purpose Changes the command prompt.
\Format \SyntaxLine{\kw{PROMPT} \opt{\ph{prompt-text}}}
\Type{Internal}
\Remarks
\Param{prompt-text} is the new prompt.  Characters are displayed as you
type them, except that a dollar sign and the character after it stand for
the following:
\par\medskip
\begin{tabular}{@{}l@{\quad}p{75mm}@{}}
{\imSideHeadFont Code} & {\imSideHeadFont Displays}\\[2pt]
\kw{\$P} & The current drive and directory, for example \texttt{A:\bs SAMPLES}\\
\kw{\$G} & The \texttt{>} character\\
\kw{\$L} & The \texttt{<} character\\
\kw{\$B} & The \texttt{|} character\\
\kw{\$Q} & The \texttt{=} character\\
\kw{\$N} & The current drive letter\\
\kw{\$T} & The time, as \texttt{hh:mm:ss.cc}\\
\kw{\$\_} & Starts a new line\\
\end{tabular}
\par\medskip
A dollar sign followed by any other character displays that character.

\draft{bead needed: in the shipped build \$D displays 00-00-0000 and \$\$
ends the prompt at that point; DOS 3.3 displays the date and a dollar sign}
\kw{\$D} displays the date, and \kw{\$\$} displays the \texttt{\$}
character.

Typed with no parameter, \cmd{PROMPT} restores the standard prompt,
\kw{\$P\$G}.  The prompt is kept in the environment variable
\texttt{PROMPT}, which \cmd{SET} displays.  If there is no room for it,
\cmd{PROMPT} displays \texttt{Out of environment space}.
\Example
The first command makes the prompt the drive letter alone; the second makes
it the word \texttt{Ready}; the third restores the standard prompt:
\Shot{prompt}
\end{DosEntry}

% RD -- os/milton/command.c builtin_rd (MSG-DOS-0011, MSG-DOS-0019; silent
% success).  Gates: test-ut6d, test-fat-subdir.
\begin{DosEntry}{RD}{Remove Directory}{internal}
\AlsoNamed{RMDIR}
\Purpose Removes a subdirectory.
\Format \SyntaxLine{\kw{RD} \opt{\ph{d}\kw{:}}\ph{path}}
  \SyntaxLine{\kw{RMDIR} \opt{\ph{d}\kw{:}}\ph{path}}
\Type{Internal}
\Remarks
\Param{path} names the directory to remove.  The directory must be empty
except for its \file{.} and \file{..} entries.  You cannot remove the root
directory or the current directory.

\cmd{RD} displays nothing when it succeeds.  If the directory does not
exist, is not empty, or is the current directory, \cmd{RD} displays
\texttt{Invalid path, not directory, or directory not empty} and nothing is
removed.  Delete the files in a directory with \cmd{DEL} before you remove
it.
\Example
The example creates and removes the directory \file{TEMP}, then tries to
remove \file{SAMPLES}, which still holds files:
\Shot{rd}
\end{DosEntry}

% REN -- os/milton/command.c builtin_ren (AH=56h, same directory; failure
% MSG-DOS-0002; MSG-DOS-0011; silent success).  Gates: test-command,
% test-fat12-dir.
\begin{DosEntry}{REN}{Rename}{internal}
\AlsoNamed{RENAME}
\Purpose Changes the name of a file.
\Format \SyntaxLine{\kw{REN} \opt{\ph{d}\kw{:}}\opt{\ph{path}}\ph{filename}\opt{\kw{.}\ph{ext}}
  \ph{filename}\opt{\kw{.}\ph{ext}}}
  \SyntaxLine{\kw{RENAME} \opt{\ph{d}\kw{:}}\opt{\ph{path}}\ph{filename}\opt{\kw{.}\ph{ext}}
  \ph{filename}\opt{\kw{.}\ph{ext}}}
\Type{Internal}
\Remarks
The first file specification names the file to rename.  The second is the
new name; do not give it a drive or path.  The file stays in its directory.
Wildcards are not allowed in this version.

\cmd{REN} displays nothing when it succeeds.  If the file does not exist,
or a file with the new name already exists, \cmd{REN} displays
\texttt{Bad command or file name}.  With only one name it displays
\texttt{Required parameter missing}.
\Example
To rename \file{MEMO.TXT} to \file{MEMO.OLD} and list the directory:
\Shot{ren}
\end{DosEntry}
